% 本文件由 LaTeX 排版 Agent 根据有效 translation.json 自动生成。
% author=Ignatius of Antioch; work_id=0110; title=依纳爵致波利卡普书

\chapter{依纳爵致波利卡普书}

% structure: p0001 source=h1 target=omitted reason=page-title-already-rendered-by-parent

% structure: p0002 source=h2 target=section reason=relative-heading-rank; base=h2

\section{问候}

依纳爵，又称德敖弗罗，致斯米纳教会的主教波利卡普——其实他拥有天主父和主耶稣基督作为自己的主教：深祝幸福。\par

% structure: p0004 source=h2 target=section reason=relative-heading-rank; base=h2

\section{第一章 称赞与劝勉}

我已确知你的心志坚定不移于天主，如同在不可动摇的磐石上，因此我高声光荣祂的圣名，使我堪当目睹你那无可指责的面容——愿我常在主内享见你！我以你所穿戴的恩宠恳求你，在赛程中竭力前进，并劝勉众人，使他们得救。你要以全部谨慎持守你的岗位，无论在肉身上还是精神上。努力维系合一，没有比这更好的事了。要容忍一切，如同主容忍你一样。在爱德中支持众人，正如你此刻所做。要不断祈祷。\BibleRef{得前 5:17} 求你已有的理解之外再加求更多的智慧。要警醒，拥有不眠的灵魂。要按天主赐给你的能力，分别对每个人说话。担负众人的软弱，如同一位完美的运动员：哪里劳苦大，收获也就越大。\par

% structure: p0006 source=h2 target=section reason=relative-heading-rank; base=h2

\section{第二章 劝勉}

如果你爱善良的门徒，这并不值得称赞；反而当以温柔去降服那些更难对付的人。并非每种伤口都可用同一膏药医治。用温和的敷料减轻剧烈的攻击。在一切事上，要\Quote{灵巧像蛇，纯良像鸽子}。\BibleRef{玛 10:16} 为此，你由肉体和精神组成，好能温柔对待那些显现在你眼前的恶。至于那些看不见的，要祈求天主向你启示它们，使你在任何恩赐上都不缺少，反而丰盈有余。时代需要你，如同舵手需要风，如同在风暴中颠簸的人寻求港口，好使你和你的羊群都能抵达天主那里。要清醒，如同天主的运动员；摆在你面前的奖赏是不朽和永生——你对此也深信不疑。在一切事上，愿我的灵魂为你，以及我所佩戴的锁链——你所爱的——而奉献。\par

% structure: p0008 source=h2 target=section reason=relative-heading-rank; base=h2

\section{第三章 劝勉}

不要让那些看似可信、却传授异端道理的人使你恐惧。\BibleRef{弟前 1:3} 要站稳，如同受捶打的铁砧。高贵的运动员被击伤却仍得胜，这是他的本分。尤其我们应为天主的缘故忍受一切，好使祂也容忍我们。你当比现在更加热切。要仔细权衡时机。期待那位超越一切时间、永恒而不可见、却为了我们而成为可见者；不可捉摸、不可受苦、却为了我们而成为可受苦者；并在各种方式中为我们受苦的那一位。\par

% structure: p0010 source=h2 target=section reason=relative-heading-rank; base=h2

\section{第四章 劝勉}

不要让寡妇被忽视。你要效法主，作她们的护卫和朋友。凡事不可未经你同意而行；你也不可未经天主批准而行——你确实不会如此，因你坚定。你们要频繁聚会：逐一提名寻找众人。不要轻视男女奴隶，但也不要让他们骄傲自大，反而要他们更加顺服，为光荣天主，好使他们从天主那里获得更美好的自由。不要让他们渴望用公费获得自由，免得成为自己私欲的奴隶。\par

% structure: p0012 source=h2 target=section reason=relative-heading-rank; base=h2

\section{第五章 夫妻的本分}

逃避邪恶的技艺，但更要公开谈论它们。对我的姐妹们说，要爱主，要在肉身和精神上满足于自己的丈夫。同样，以耶稣基督的名劝勉我的弟兄们，要爱自己的妻子，如同主爱教会。\BibleRef{弗 5:25} 如果有人能保持贞洁，以光荣那掌管肉身的主，就让他不夸耀地保持。如果他开始夸耀，他就灭亡了；如果他自认比主教还大，他就毁了。但无论男女，若结婚，就应在主教赞同下结合，使他们的婚姻合乎天主，而非顺从自己的私欲。一切事都要为光荣天主而行。\BibleRef{格前 10:31}\par

% structure: p0014 source=h2 target=section reason=relative-heading-rank; base=h2

\section{第六章 基督徒羊群的本分}

要留心主教，好使天主也留心你。愿我的灵魂为那些顺服主教、长老和执事的人，但愿我与他们一同在天主里！你们要彼此合作，一同努力，一同奔跑，一同受苦，一同睡觉，一同醒来，如同天主的管家、同工和仆人。要取悦你们为之争战并从祂那里领受酬劳的那位。不要让你们中任何人被发现是逃兵。愿你们的圣洗圣事作为武器，信德作为头盔，爱德作为长矛，忍耐作为全副武装。愿你们的善工成为交付给你们的任务，好使你们领受相称的回报。因此，要彼此宽容，存心温柔，如同天主对待你们一样。愿我永远从你们获得喜乐！\par

% structure: p0016 source=h2 target=section reason=relative-heading-rank; base=h2

\section{第七章 请求波利卡普派遣使者往安提约基雅}

据报告得知，叙利亚安提约基雅的教会因你们的祈祷而平安，我也就更加受到鼓励，在无忧无虑中安息于天主内——若我确能借苦难抵达天主，那么通过你们的祈祷，我就能被算作基督的门徒。最蒙福于天主的波利卡普，应当召集一个极其庄重的会议，选出一位你深爱、且知道是活跃之人，他被指定为天主的使者；并赐予他这份殊荣，使他前往叙利亚，并光荣你那常在的爱德，为基督的赞美。基督徒不拥有自己的主权，而必须时常为天主的服务准备就绪。现在，这工作既是天主的，也是你的，当你为祂的光荣完成它时。因为我深信，你已借着恩宠预备好作一切属于天主的善工。因此，知道你热切爱慕真理，我以此简短书信劝勉你。\par

% structure: p0018 source=h2 target=section reason=relative-heading-rank; base=h2

\section{第八章 也让其他教会派使者往安提约基雅}

因为我不曾能够写信给所有的教会，因为必须突然从\Place{特洛阿}起航去\Place{乃阿颇里}，如[皇帝的]旨意所命，[我恳求你]，你既熟知天主的旨意，将写信给邻近的教会，叫他们也照样行事，凡是能做的就派使者，其余的则通过你所派遣的人传递信件，好使你因这永受纪念的工作而受光荣，正如你确所堪当的。我一一提名问安，特别问候\Person{埃庇特洛普}的妻子，以及她的全家与子女。我问候我所爱的\Person{阿塔路}。我问候那被认为堪当[从你们那里]往\Place{叙利亚}去的人。恩宠将永远与他同在，也与那派遣他的\Person{波利卡普}同在。我祈求你们的幸福永远在我们的天主耶稣基督内，藉着他，你们继续在合一与天主的保护之下。我问候我亲爱的\Person{阿尔刻}。在主内安好。\par

\SourceRecord{Translated by Alexander Roberts and James Donaldson.；Ante-Nicene Fathers；Vol. 1.；Edited by Alexander Roberts, James Donaldson, and A. Cleveland Coxe.；Buffalo, NY: Christian Literature Publishing Co.,；1885.；<http://www.newadvent.org/fathers/0110.htm>.}
